A relação dimensional da eq. (3.10) é verificada expressando cada grandeza em dimensões fundamentais $[M,\,L,\,T]$ e confirmando a homogeneidade.

\textbf{Dimensões das grandezas envolvidas:}
\begin{itemize}
\item Potência mecânica: $[P] = M\,L^{2}\,T^{-3}$
\item Força: $[F] = M\,L\,T^{-2}$
\item Velocidade: $[v] = L\,T^{-1}$
\item Energia cinética: $[E_k] = M\,L^{2}\,T^{-2}$
\item Taxa de variação de energia: $\bigl[dE_k/dt\bigr] = M\,L^{2}\,T^{-3}$
\end{itemize}

\textbf{Verificação explícita:} O produto $F\cdot v$ tem dimensão:
\[
[F\cdot v] = [M\,L\,T^{-2}]\cdot[L\,T^{-1}] = M\,L^{2}\,T^{-3} = [P].\quad\checkmark
\]
Igualmente, $dE_k/dt$ tem dimensão:
\[
\left[\frac{dE_k}{dt}\right] = \frac{M\,L^{2}\,T^{-2}}{T} = M\,L^{2}\,T^{-3} = [P].\quad\checkmark
\]
Portanto, ambos os lados da eq. (3.10) possuem a mesma combinação de $M$, $L$ e $T$, confirmando que a relação dimensional é correta e a equação é fisicamente consistente.

